\lecture{11}{Tue 13 Feb 2024 13:44}{}

\begin{theorem}
	\label{thm:SubdivisionNatural}
	
	$\\\$ $ is a natural chain map and is naturally chain homotopic to Id via $T$.
\end{theorem}
\begin{proof}(First Claim)
	Let $\sigma : \Delta_n \to Y$ be affine. The $n = 0$ case is simple. Assume true for $n -1$. Then by affine-ness, $f(b_\sigma) = b_{f_*(\sigma)}$. 

	\begin{align*}
		f_*(\$ (\sigma)) 
		&= f_*\left( b_\sigma * \$ d \sigma \right)  \\
		&= f(b_\sigma) * f_*(\$ d\sigma) \\
		&= b_{f_*(\sigma)} * f_*(\$ d\sigma) \\
		&= b_{f_*(\sigma)} * \$ f_*(d \sigma)
	.\end{align*}
	On the other hand,
	\[
		\$ f_*(\sigma) = b_{f_*(\sigma)} * \$ d f_*(\sigma)
	.\] 

	Since $f_*$ is a chain map, the two sides are equal.
\end{proof}

Now, we want to extend $\$$ to arbitrary simplices in arbitrary topological spaces. Let $X$ be a topological space. Pick $\sigma \in \text{Sin}_n(X)$. We first consider affine simplices in the standard simplices, and then push forward to singular simplices:

% https://quiver.local:8080/#q=WzAsNyxbMSwwLCJcXERlbHRhX24iXSxbMiwwLCJYIl0sWzIsMiwiU19uKFxcRGVsdGFfbikiXSxbNCwyLCJTX24oWCkiLFszMCw2MCw2MCwxXV0sWzAsMSwiU19uXlxcdGV4dHthZmZ9KFxcRGVsdGFfbikiLFswLDYwLDYwLDFdXSxbMCwwLCJcXERlbHRhX24iXSxbMCwyLCJTX25eXFx0ZXh0e2FmZn0oXFxEZWx0YV9uKSJdLFswLDEsIlxcc2lnbWEiLDAseyJjb2xvdXIiOlszMCw2MCw2MF19LFszMCw2MCw2MCwxXV0sWzIsMywiXFxzaWdtYV8qIl0sWzUsMCwiXFxpb3RhIiwwLHsiY29sb3VyIjpbMCw2MCw2MF19LFswLDYwLDYwLDFdXSxbNCw2LCJcXFxcXFwkIiwyXSxbNiwyLCIiLDAseyJzdHlsZSI6eyJ0YWlsIjp7Im5hbWUiOiJob29rIiwic2lkZSI6InRvcCJ9fX1dXQ==
\[\begin{tikzcd}
	{\Delta_n} & {\Delta_n} & X \\
	\textcolor{rgb,255:red,214;green,92;blue,92}{S_n^\text{aff}(\Delta_n)} \\
	{S_n^\text{aff}(\Delta_n)} && {S_n(\Delta_n)} && \textcolor{rgb,255:red,214;green,153;blue,92}{S_n(X)}
	\arrow["\sigma", color={rgb,255:red,214;green,153;blue,92}, from=1-2, to=1-3]
	\arrow["{\sigma_*}", from=3-3, to=3-5]
	\arrow["\iota", color={rgb,255:red,214;green,92;blue,92}, from=1-1, to=1-2]
	\arrow["{\\\$}"', from=2-1, to=3-1]
	\arrow[hook, from=3-1, to=3-3]
\end{tikzcd}\]

\begin{example}
	Take $n = 1$, consider some $\sigma: \Delta_1 \to X$. Let $v = b i_n = \frac{1}{2}(e_0 + e_1)$. Then $\$(i_n) = [v, e_1] - [v, e_0]$. Thus, $\$ \sigma = \sigma_*([v, e_1 - [v, e_0] ) = \sigma|_{[v, e_1]} - \sigma|_{[v, e_0]}$.
\end{example}

\begin{proof}(Chain Homotopy)
	\begin{align*}
		d \$ (\sigma)
		&= d \sigma_* \$(i_n) \\
		&= \sigma_* d \$ (i_n) \\
		&= \sigma_* \$ d(i_n) 
	.\end{align*}
	Let $d^i: \Delta_{n  - 1} \to  \Delta_n$ be affine face-inclusion maps. Then
	\begin{align*}
		\$ d(\sigma) &= \$ \left( \sum_{i = 0}^n (-1)^i \sigma \circ d^i \right) \\
		&= \sum_i (-1)^i \left( \sigma \circ d^i)_* \$(i_{n - 1}) \right)  \\
		&= \sigma_* \left( \sum_i (-1)^i (d^i)_* \$ i_{n - 1} \right)  \\
		&= \sigma_* \left( \sum_i (-1)^i  \$ (d^i)_* i_{n - 1} \right)  \\
		&= \sigma_*  \$ \left( \sum_i (-1)^i  (d^i)_* i_{n - 1} \right)  \\
		&= \sigma_*  \$ d(i_n)
	.\end{align*}

Currently we only have $T $ for the affine case; we now extend it to the general case.
\begin{align*}
	T: S_n(X) &\longrightarrow S_{n + 1}(X) \\
	\sigma &\longmapsto T(\sigma) = \sigma_* T(i_n)
.\end{align*}

For $\sigma: \Delta_n \to X$,
\begin{align*}
	d T \sigma + T d \sigma
	&= d \sigma_* T(i_n) + T\left( \sigma_i (-1)^i \sigma \circ d^i \right) \\
	&= \sigma_* d T(i_n) + \sum_i (-1)^i \sigma_* (d^i)_* T(i_{n - 1}) \\
	&= \sigma_* d T(i_n) + \sum_i (-1)^i \sigma_* T (d^i)_* (i_{n - 1}) \\
	&= \sigma_* d T(i_n) + \sigma_* T \sum_i (-1)^i (d^i)_* (i_{n - 1}) \\
	&= \sigma_* d T(i_n) + \sigma_* T d i_n \\
	&= \sigma_* (d T(i_n) + T d i_n) \\
	&= \sigma_* (i_n - \$ i_n) \\
	&= \sigma - \$ \sigma 
.\end{align*}
\end{proof}

\begin{corollary}
	\label{cor:Subdivisioniterate}
	\[
		\$^k \cong \text{Id}
	.\] 
\end{corollary}
\begin{proof}
	The point is that we have an explicit chain homotopy. Define
	\[
		T^k = T \left( 1 + \$ + \$^2 + \ldots + \$^{k - 1} \right) : S_n(X) \to S_{n + 1}(X)
	.\] 
	\todo{verify the corollary.}
\end{proof}

\begin{lemma}
	\label{lem:DiameterVertice}

	For  $A \subset \mathbb{R}^k$ define the diameter of $A$ to be $\text{diam}(A) = \sup_{x, y \in A} \|x - y\|$.
	
	For any $\sigma \in S_n^{\text{aff}}(X)$, $\text{diam}(\sigma) := \text{diam}(\text{Im}(\sigma)) $ is achieved at a pair of vertices of $\sigma$.
\end{lemma}
\begin{proof}
	Let $\sigma_0 = [v_0, \ldots, v_n]$, let $v, w \in \operatorname{Im} \sigma$. Then if $t_i \ge 0, \sum t_i = 1$,
	\begin{align*}
		\|v - w\| 
		&= \| v - \sum_{i = 0}^n t_i v_i\|\\
		&= \| \sum_{i = 0}^n t_i v - \sum_{i = 0}^n t_i v_i\|\\
		&= \| \sum_{i = 0}^n t_i (v -  v_i)\|\\
		&\leq \sum_{i = 0}^n t_i \|v -  v_i\|\\
		&\leq  \max_{i} \|v - v_i\|
	.\end{align*}

	Apply this argument to itself, we get
	\[
		\|v - v_i\| \le \max_j \|v_j - v_i\|
	.\] 
	Therefore
	\[
		\| v- w\| \le \max_{i, j} \|v_j - v_i\|
	.\] 
\end{proof}

\begin{lemma}
	\label{lem:SubdivisionDiameter}
	Let $\sigma \in S_n^{\text{aff}}(Y)$ and $Y \subset \mathbb{R}^N$. Let $\tau$ be a simplex in the sum of $\$ \sigma$. Then
	\[
		\text{diam} (\tau) \le  \frac{n}{n + 1} \text{diam} (\sigma)
	.\] 
\end{lemma}
\begin{proof}
	Let $b_1, b_2$ be two vertices of $\tau$. Then there exists sub-simplices $f_1 \subset f_2 \subset \sigma$, $b_1 = b_{f_1}, b_2 = b_{f_2}$. There can be dimension gap, so they are not in general faces. 
	Without loss of generality, assume $f_2 = [v_0, \ldots, v_k], f_1 = [v_0, \ldots, v_n]$, where $m < k \le n$. Now
	\[
		b_1 = \frac{1}{m + 1} \sum_{i = 0}^m v_i 
		\qquad
		b_2 = \frac{1}{k + 1} \sum_{i = 0}^k v_i
	.\] 
	\begin{align*}
		\|b_2 - b_1\| = \|\sum_{i = 0}^m \frac{1}{m + 1} (b_2 - v_i)\| \le  \max_{0 \le  i \le  m} \|b_2 - v_i\| \le  \max_{0 \le  i \le  k} \|b_2 - v_i\|
	.\end{align*}
	WLOG, assume $\|b_2 - v_0\| = \max_{0 \le  i \le  k} \|b_2 - v_i\|$. Now
	\begin{align*}
		\|b_2 - v_0\| 
		&= \|\frac{1}{k + 1} \sum_{i = 0}^k v_i - v_0\|\\
		&= \frac{k}{k + 1} \|\frac{1}{k} \sum_{i = 1}^k v_i - v_0\| \\
		&\leq \frac{k}{k + 1} \max_{1 \le i \le k}\| v_i - v_0\| \\
		&\leq \frac{n}{ n + 1} \text{diam}(\sigma)
	.\end{align*}
\end{proof}

